\documentclass[quick,con]{debate}
\motion{“磕CP”热潮折射出人们对爱情的更深信仰 / 更深怀疑}
\stance{反方}{更深怀疑}
\begin{document}
\debatetitle
\begin{thesisbox}人们只敢把爱放在屏幕那头、结局写好的地方；信的是别人的爱，疑的是自己的爱靠不靠得住，疑在底下。\end{thesisbox}

\section{定义与判准（标准表述）}
\begin{fields}
\field{嗑CP}{以旁观者的身份，给虚构角色或者真实人物配对，想象并期待他们之间的浪漫关系，并从中得到情感满足。}
\field{信仰 / 怀疑}{信仰：相信真正的爱情存在、值得期待和投入。怀疑：不相信它存在，或者不相信它靠得住、值得投入。}
\field{更深}{两种心态都有，比哪一种在底下，也就是谁是另一种的来源。}
\field{判准}{哪一种心态更能解释热潮的核心特征：嗑的是什么样的爱，为什么用旁观的方式去爱；并且是另一种心态的来源。\sub{我方要证明}{嗑CP 的动力来自不信真爱靠得住；热潮里的甜只在安全距离内成立。}\sub{对方要证明}{嗑CP 的动力来自相信真爱存在。}}
\field{对方提偏向定义或判准时的 30 秒口径}{对方说“喜欢爱情故事就算信”：那看言情小说的人全都信爱情，题目就没有意义了。中立定义里的怀疑包括不信它靠得住、值得投入，请回到判准的另一半：为什么只敢这样爱。}
\end{fields}

\section{我方论点}
\begin{tabularx}{\linewidth}{@{}L{4.6em}Y Y L{6.4em}@{}}\toprule
\thead{标签}&\thead{一句话}&\thead{核心数据 / 学理 / 案例}&\thead{出处}\\\midrule
\textbf{信在远处}&爱永远在屏幕那头；人们信爱会发生在别人身上，却不信、也不敢让它落到自己身上&数据：2905 名未婚城市青年，12.8\% 将来不谈，26.3\% 不确定；案例：“虽然我不谈恋爱，但是我要看我的CP谈恋爱”；二辩新一层：观众几乎不用付出，随时可以抽身&共青团中央课题组 2021；许冠文 2023；Horton 和 Wohl 1956\\
\textbf{要的是保险}&结局自己写好、不会解散、不会受伤；只有结局写好才敢投入，说明不信真实的爱靠得住&案例：受访者“在现实里体验这种爱情，自己很难不发疯”；许冠文“你已经给他们画好了终点”；三辩新一层：巴迪欧“安全威胁”；不想恋爱的人里 74.8\% 选“一个人很好，谈恋爱很麻烦”&许冠文 2023；Badiou 2009；共青团中央课题组 2021\\
\bottomrule\end{tabularx}

\section{对方论点 → 一句话反驳}
\begin{tabularx}{\linewidth}{@{}Y Y Y@{}}\toprule
\thead{对方会说}&\thead{我方一句话回}&\thead{对方追问时}\\\midrule
嗑的是理想：专一、双向，所以信&标准高是因为不用兑现；叶公对龙的标准也很高&“标准带回现实：宁缺毋滥”→ 落到现实就是缺\\
疑的是条件：76.8\% 想谈，卡在圈子&想要不等于相信；想而不敢正是怀疑&“疑遇见不是疑爱情”→ 扛不住圈子窄的爱，在人心里就是靠不住\\
怕风险说明把爱当真&当真和相信它靠得住是两件事；当真又不敢，就是定义里的怀疑&“那你们承认有信”→ 承认，信在上面，底下是不敢\\
许冠文说“赋予了很高的价值”&同一句后半：“每一步都充满风险”，决定人们怎么做的是后半&“前半才是底”→ 那为什么只敢放在屏幕那头？\\
隔着屏幕的关系对多数人是补充&补充的前提是正餐还在吃；一成多说不谈，超过四分之一不确定&“那是全体青年”→ 背景数据；主证据是嗑的人自己的话\\
\bottomrule\end{tabularx}

\section{必问与必答}
\begin{points}{必问（对方不答就点名、重复、不换题）}
\item 如果人们真信爱情，为什么要把它放在一段自己碰不到、结局早写好的关系里？（全场主问题）
\item 既然信，为什么不去谈？
\item 不敢把爱放在自己身上，是不是不信它靠得住？
\end{points}
\begin{points}{必答（15 秒正面回答）}
\item 不信爱情的人为什么嗑专一？ → 只有放在不会伤到自己的关系里，才敢要专一。
\item 旁观就是怀疑？ → 不是所有旁观，是说出“现实太险”的旁观。
\item 234 份问卷能代表谁？ → 不能代表全体，我方只用它说明有这样的心态，主证据是许冠文的研究。
\end{points}

\section{对辩战场概要}
\begin{tabularx}{\linewidth}{@{}L{5.2em}Y Y Y@{}}\toprule
\thead{对辩}&\thead{开场问题}&\thead{目标}&\thead{收束语}\\\midrule
三辩（接第一战场）&专一为什么只敢放在进不去的关系里？&把内容拉到形式：方式说明敢不敢&内容说明想要什么，方式说明敢不敢\\
一辩（收拢）&是因为不敢才去嗑，对吗？&来源：不敢在底下&先有不敢，才去屏幕那头找\\
二辩（结辩前）&判准的两半，对方只讲了哪一半？&守判准的“为什么这样嗑”&信在上面，疑在底下\\
\bottomrule\end{tabularx}

\section{如果……就……}
\begin{contingency}
\ifthen{对方把信仰定义成“喜欢爱情故事就算”}{说出 30 秒口径，拉回判准的另一半。}
\ifthen{对方拿出没预料到的数据}{先问出处，再问它测的是想要，还是相信它靠得住。}
\ifthen{论点一被打成“旁观不等于怀疑”}{收缩到论点二：结局写好才敢投入，是对靠不靠得住的判断。}
\end{contingency}
提醒：反二质询正一是全场第一次质询，为反一立论开头取材；反二申论是全场第一篇；反三申论先回应正三；反一先结辩，只能封路。质询别频繁打断，被质询先答是或否；对方回避主问题就点名、重复、不换题。

\begin{recordcard}
\record{反二质询后}{质询拿到了什么}{\opt{质询拿到了正一承认“嗑CP的人不进入那段关系”}\opt{质询里正一回避了“既然信，为什么不去谈”}\opt{质询没有拿到明确成果}}{反一开篇立论·开头}
\record{正一立论后}{正方立论的主干}{\opt{正方立论的主干是“嗑的是理想，所以是信”}\opt{正方立论的主干是“疑的只是遇见，不是爱情”}\opt{正方立论的主干不在以上两条时}}{反二申论·拆正方立论}
\record{正三申论后}{正三申论主打}{\opt{正三申论主打“婚、育、恋里，恋爱被放下得最少”}\opt{正三申论主打“对多数人是补充，不是替代”}\opt{正三申论主打不在以上两条时}}{反三申论·回应正三申论}
\record{正三申论后}{正方到此最有力的一点}{\opt{正方全场最有力的是“嗑的是专一、双向，不信的人不会嗑这些”}\opt{正方全场最有力的是“怕风险，说明把爱当真”}\opt{正方最有力的一点不在以上两条时}}{反三申论·处理正方最强点}
\record{二辩对辩后}{全场主战场}{\opt{全场主战场落在“嗑的内容说明什么”}\opt{全场主战场落在“不敢算不算不信”}\opt{主战场不在以上两条时}}{反一结辩·归纳分歧}
\record{二辩对辩后}{对方全场最有力的一点}{\opt{正方全场最有力的是“嗑的是专一和双向”}\opt{正方全场最有力的是“疑的只是遇见”}\opt{正方最有力的一点不在以上两条时}}{反一结辩·处理最强点}
\record{二辩对辩后}{正方全场主打的路线}{\opt{正方全场主打“内容是理想，所以是信”}\opt{正方全场主打“疑的是门口，门里有人”}\opt{正方全场路线不在以上两条时}}{反一结辩·封路}
\record{全场}{对方回避了哪些问题}{（写下来）}{申论与结辩里点名}
\end{recordcard}
\end{document}
